% Lecture 09: Templates and generic library design
\section{Лекция 9. Templates: generic code как библиотека}
\label{sec:lecture-09}

В блоке~1 запись \texttt{std::vector<int>} читалась как «vector элементов
int», а в блоке~3 появилась \texttt{std::array<int, 4>}. Теперь разберём
механизм языка, благодаря которому библиотеке не нужны вручную написанные
\texttt{vector\_int} и \texttt{vector\_string}. Главный вопрос блока:
как один исходный алгоритм или тип сделать reusable для разных типов, не
стирая их различия во время выполнения?

\begin{importantbox}[Причинная цепочка]
Дублирование по типам $\rightarrow$ function template $\rightarrow$
deduction $\rightarrow$ instantiation $\rightarrow$ requirements к типу
$\rightarrow$ class template и value parameter $\rightarrow$ маленькая
generic library $\rightarrow$ мост к STL.
\end{importantbox}

\subsection{Design pressure: алгоритм тот же, тип другой}

До templates максимум двух типов пришлось бы записать дважды:

\begin{cppcode}
int max_int(int left, int right) {
    return left < right ? right : left;
}

double max_double(double left, double right) {
    return left < right ? right : left;
}
\end{cppcode}

Имена и типы различаются, но решение одно: сравнить два значения и вернуть
большее. Исправление ошибки или изменение правила придётся повторять. Overload
делает единое имя вызова, но тела всё ещё дублируются. Common base и virtual
function из блока~7 здесь искусственны: \texttt{int} и \texttt{double} не
должны становиться объектами одной hierarchy.

\term{Generic programming} описывает алгоритм через необходимые операции, а
не через один заранее выбранный concrete type. Языковой механизм templates
позволяет compiler получить concrete entities для подходящих arguments.

\subsection{Function template по частям}

\begin{cppcode}
template <typename T>
const T& maximum(const T& left, const T& right) {
    return left < right ? right : left;
}
\end{cppcode}

\texttt{template <typename T>} вводит \term{template parameter} \texttt{T}:
здесь это parameter-типа. В этом месте \texttt{class T} синтаксически тоже
допустимо; в конспекте последовательно используем \texttt{typename}, потому
что parameter может быть не class, а, например, \texttt{int}.

Всё определение --- \term{function template declaration}. Это не одна
обычная функция, способная runtime принимать значения «без типа». Template
задаёт семейство возможных functions. Записи \texttt{maximum<int>} и
\texttt{maximum<double>} называют разные \term{specializations} этого
template с разными function types.

\begin{standardbox}[Template, specialization и instantiation]
Template --- исходное параметризованное объявление. Specialization --- entity,
полученная для конкретных template arguments или заданная отдельно.
Instantiation --- процесс получения specialization из template. В простом
вызове compiler не обязан буквально копировать текст; это лишь удобная первая
модель, а не требование к реализации или object code.
\end{standardbox}

\subsection{Template argument deduction}

Обычно argument типа выводится из function arguments:

\begin{cppcode}
int a = maximum(4, 9);              // T deduced as int
double b = maximum(2.0, 3.5);       // T deduced as double
std::string c = maximum(std::string{"a"}, std::string{"b"});
\end{cppcode}

Это \term{template argument deduction}: compiler сопоставляет parameter types
\texttt{const T\&} с types expressions в call и выводит \texttt{T}. Вызов не
делает type runtime-данными; решение принимается при translation.

Template argument можно указать явно:

\begin{cppcode}
double value = maximum<double>(2, 3.5);
\end{cppcode}

Здесь первый argument \texttt{int} преобразуется к явно выбранному
\texttt{double}. Явная запись полезна, когда deduction невозможен или нужен
осознанный target type, но очевидный argument обычно не повторяют.

Если оба parameters содержат один \texttt{T}, compiler должен получить
согласованный результат. Вызов \texttt{maximum(2, 3.5)} пытается вывести из
первого argument \texttt{int}, а из второго \texttt{double}; простое правило
не выбирает «лучший общий тип», и call ill-formed. Можно привести arguments к
одному смысловому типу или явно выбрать \texttt{maximum<double>}.

Это только базовая deduction. Сложные правила, \texttt{auto}, forwarding и
class template argument deduction относятся к следующим блокам.

\subsection{Instantiation: что требуется compiler}

Когда program использует \texttt{maximum(4, 9)}, происходит implicit
instantiation подходящей specialization \texttt{maximum<int>}. Использование
со strings создаёт \texttt{maximum<std::string>}. Каждая проверяется как
concrete typed function.

Стандарт задаёт правила well-formed program и observable behavior, но не
обещает отдельную физическую копию machine code для каждой specialization.
Compiler/linker могут inline, объединять эквивалентный code или вообще не
оставить callable symbol после optimization. Обратное тоже возможно: разные
specializations дают разные инструкции и занимают место.

\subsection{Requirements: generic не означает «любой тип»}

Implementation \texttt{maximum} содержит \texttt{left < right}. Значит,
выбранный \texttt{T} должен предоставлять подходящую операцию \texttt{<}, а её
результат должен работать как condition. Это \term{requirement} алгоритма.

\begin{importantbox}[Обязательная формулировка]
Template не «принимает любой тип». Он работает с теми типами, которые
удовлетворяют операциям, используемым implementation.
\end{importantbox}

Для следующей функции requirements другие:

\begin{cppcode}
template <typename Container, typename Value>
std::size_t count_matching(const Container& values, const Value& wanted) {
    std::size_t count{};
    for (const auto& value : values) {
        if (value == wanted) {
            ++count;
        }
    }
    return count;
}
\end{cppcode}

\texttt{Container} должен поддерживать range-for, а его element --- сравнение
с \texttt{Value} через \texttt{==}. Функция не требует именно vector и не
обещает весь interface контейнера. Range-for здесь использует уже знакомую
форму; полный iterator model будет в блоке~14.

В простом C++17 interface эти requirements явно не записаны. Они проверяются
при instantiation. Если вызвать \texttt{maximum} для \texttt{NotOrdered} без
\texttt{operator<}, compiler укажет ошибку внутри template и цепочку вроде
«required from here». Это compile-time error, а не runtime fallback.

\subsection{Как читать template diagnostics}

Длинная diagnostic часто содержит три уровня:
1) пользовательский call, который запросил instantiation; 2) concrete template arguments, например \texttt{T = NotOrdered}; 3) первая неподходящая operation внутри definition.

Формат GCC и Clang различается. Практический порядок: найти первый свой call,
установленные template arguments и ближайшую первичную причину, а не исправлять
все последующие сообщения. Negative tests сохраняют настоящие diagnostics в
build-каталог, но PDF не копирует многоэкранный compiler dump.

В C++20 Concepts позволяют выразить часть requirements в interface и улучшить
диагностику. Здесь мы не вводим их синтаксис: сначала важно увидеть сам
контракт операций.

\subsection{Value, const reference и mutable reference}

Templates не отменяют семантику блока~4:
\texttt{T value} создаёт отдельный parameter object; это уместно для маленького значения или нужной копии; \texttt{const T\& value} читает существующий object без намеренной копии и подходит \texttt{maximum}; \texttt{T\& value} означает изменение caller's object.

Выбор определяется contract, lifetime и cost, а не словом \texttt{template}.
\texttt{maximum} возвращает \texttt{const T\&}; в реальной библиотеке caller
не должен сохранять эту reference дольше выбранного argument. Для простых
вычислений часто безопаснее вернуть \texttt{T} по value. Forwarding references
и perfect forwarding здесь намеренно не нужны.

Простой return type \texttt{T} понятен, когда operation действительно
возвращает тот же type. Смешанное сложение разных types потребовало бы иной
модели результата; не будем преждевременно вводить \texttt{decltype} machinery.

\subsection{Class template: семейство concrete types}

Если параметризуется не одна operation, а состояние вместе с interface,
нужен class template:

\begin{cppcode}
template <typename T>
class Box {
public:
    explicit Box(const T& value) : value_{value} {}
    const T& value() const { return value_; }
private:
    T value_;
};

Box<int> count{9};
Box<std::string> title{"templates"};
\end{cppcode}

\texttt{Box} --- template, а \texttt{Box<int>} и
\texttt{Box<std::string>} --- разные concrete class types. Objects \texttt{count}
и \texttt{title} имеют соответствующие specializations. Между
\texttt{Box<int>} и \texttt{Box<double>} автоматически нет inheritance
relation и implicit runtime substitutability.

Methods, определённые внутри class template, используют тот же \texttt{T}.
При внешнем definition параметр и specialization класса нужно повторить:

\begin{cppcode}
template <typename T>
const T& Box<T>::value() const {
    return value_;
}
\end{cppcode}

Для коротких methods inline definition обычно читается проще.

\subsection{Non-type template parameter}

Template parameter может быть compile-time value:

\begin{cppcode}
template <typename T, std::size_t N>
class FixedBuffer {
public:
    constexpr std::size_t size() const { return N; }
    T& operator[](std::size_t index) { return values_[index]; }
private:
    std::array<T, N> values_{};
};
\end{cppcode}

\texttt{T} --- type template parameter, \texttt{N} --- non-type template
parameter типа \texttt{std::size\_t}. Значение \texttt{N} известно при
translation и входит в type: \texttt{FixedBuffer<int, 3>} и
\texttt{FixedBuffer<int, 5>} --- разные types. Это та же основная идея, что у
знакомого \texttt{std::array<T, N>}.

В учебном \texttt{FixedBuffer} \texttt{operator[]} имеет precondition
\texttt{index < N}; runtime bounds check не добавлен. Storage принадлежит
object через \texttt{std::array}, поэтому отдельного dynamic allocation нет.
Это иллюстрация, а не попытка заменить production \texttt{std::array}.

Class templates могут иметь default arguments, например
\texttt{template <typename T = int> class Box}. Тогда \texttt{Box<>} выбирает
\texttt{int}. Default полезен только при естественном частом выборе; скрывать
важный design parameter ради короткой записи не следует.

\subsection{Primary template, specialization и overloading}

Обычное параметризованное определение называется \term{primary template}.
Иногда один type требует действительно другой implementation:

\begin{cppcode}
template <typename T>
struct ValueLabel {
    static const char* text() { return "value"; }
};

template <>
struct ValueLabel<bool> {
    static const char* text() { return "logical value"; }
};
\end{cppcode}

Вторая запись --- explicit, или full, specialization для \texttt{bool}.
Specialization повышает coupling к concrete type, поэтому её применяют редко
и осмысленно, когда общий contract сохраняется. Partial specialization class
templates существует, но нужна более поздним generic designs. Function
templates частично специализировать нельзя.

Functions чаще расширяют overloading: рядом могут быть ordinary function и
function template либо несколько templates с различимыми parameter lists.
Overload участвует в выборе подходящей function; specialization задаёт вариант
уже существующего template. Подробные правила выбора и SFINAE здесь не нужны.

\subsection{Почему definitions templates часто находятся в headers}

Для non-template function привычна схема: declaration в \texttt{.hpp}, body в
одном \texttt{.cpp}, linker находит готовое definition. Для implicit
instantiation compiler обычно должен видеть template definition и concrete
arguments в точке, где specialization требуется. Одной declaration часто
недостаточно, чтобы сгенерировать нужную specialization.

Поэтому небольшие template libraries обычно помещают полные definitions в
headers с include guard. Это не означает «template невозможно определить в
\texttt{.cpp}». Можно управлять набором через explicit instantiation, но тогда
нужно заранее перечислить поддержанные arguments; такой способ менее generic
и не относится к основной модели этой лекции.

Каждая translation unit может увидеть один и тот же header. ODR и toolchain
обеспечивают корректное объединение допустимых одинаковых template entities;
точные object symbols зависят от compiler/linker.

\subsection{Compile-time и runtime polymorphism}

В блоках~7--8 runtime Strategy использовала common abstract base и virtual
call. Caller мог выбрать concrete strategy по runtime argument, а engine
работал через base reference.

Generic algorithm использует иной extension point:

\begin{cppcode}
template <typename Blocks, typename Policy, typename Reporter>
AuditSummary run_audit(const Blocks& blocks,
                       const Policy& policy,
                       Reporter& reporter) {
    AuditSummary summary;
    for (const auto& block : blocks) {
        const AuditResult result = policy.evaluate(block);
        ++summary.checked;
        if (result.passed) {
            ++summary.passed;
        }
        reporter.on_result(result);
    }
    reporter.on_finish(summary);
    return summary;
}
\end{cppcode}

Concrete \texttt{Policy} и \texttt{Reporter} известны при instantiation. Им
не нужна общая base class: достаточно \texttt{evaluate}, \texttt{on\_result}
и \texttt{on\_finish} с подходящим смыслом. Это compile-time polymorphism.

\begin{center}
\begin{tabular}{p{0.25\textwidth}p{0.32\textwidth}p{0.32\textwidth}}
\toprule
 & Templates & Virtual interface \\
\midrule
Выбор варианта & при translation & может быть runtime \\
Common base & не обязателен & задаёт единый interface \\
Dispatch & concrete call, часто inline-able & обычно indirect virtual call \\
Расширение & новая instantiation & новый subtype через base \\
Основная цена & compile time, diagnostics, code size & interface/dispatch, lifetime design \\
\bottomrule
\end{tabular}
\end{center}

Ни один подход не всегда лучше. Runtime Strategy нужна, когда policy выбирают
из config во время выполнения, хранят разнородные strategies вместе или
подключают через stable ABI. Template policy удобна, когда concrete type
известен при сборке, common hierarchy не несёт смысла, а operations образуют
достаточный compile-time contract.

\subsection{Cost model без лозунга zero-cost}

Template abstraction не обязана добавлять runtime dispatch или хранить type
metadata в object. Compiler видит concrete type и может inline и
специализировать code. Это возможность оптимизации, а не гарантия «быстрее».

Цена переносится и в toolchain: parsing и instantiation увеличивают compile
time; длинные chains ухудшают diagnostics; specializations для разных types
могут увеличить binary size. Оптимизатор может убрать или объединить часть
code, поэтому code bloat измеряют для конкретной сборки. Benchmark и assembly
в блоке N/A: архитектурное различие не требует выдуманного доказательства
скорости.

\subsection{Эксперимент 1: mini\_generic\_library}

Практический каталог \texttt{mini\_generic\_library} содержит настоящий
header \texttt{mini\_generic.hpp}. В нём находятся:
\texttt{maximum<T>} с requirement \texttt{operator<}; \texttt{count\_matching<Container, Value>} с range-for и \texttt{operator==}; \texttt{FixedBuffer<T, N>} с type и non-type parameters; небольшая full specialization \texttt{ValueLabel<bool>}.

Positive run использует deduced \texttt{maximum<int>}, explicit
\texttt{maximum<double>}, \texttt{std::string} и честный project type
\texttt{Revision} с ordering по номеру. Подсчёт проверен на
\texttt{std::vector<int>} и \texttt{std::vector<std::string>}; buffers --- на
двух \texttt{T} и двух \texttt{N}.

\begin{shellcode}
cd examples/09_templates/mini_generic_library
cmake -S . -B build -DCMAKE_BUILD_TYPE=Debug
cmake --build build
./build/mini_generic_library
ctest --test-dir build --output-on-failure
\end{shellcode}

\begin{outputcode}
maximum<int>=9
maximum<double>=3.5
maximum<string>=beta
maximum<Revision>=9
matching<int>=2
matching<string>=2
FixedBuffer<int,3>.size=3
FixedBuffer<string,2>.size=2
specialization<bool>=logical value
MINI_GENERIC_LIBRARY=PASS
\end{outputcode}

CTest отдельно компилирует \texttt{negative/no\_less.cpp}. Type
\texttt{NotOrdered} не предоставляет \texttt{<}; PASS означает ненулевой exit
code compiler и непустую реальную diagnostic в
\texttt{build/negative/no\_less.log}, а не поддельный текст ошибки.

\subsection{Эксперимент 2: generic\_course\_audit}

Второй проект read-only инспектирует настоящий набор из 15 blocks. Один
\texttt{run\_audit} инстанцируется для \texttt{CompactReporter} и
\texttt{DetailedReporter}. Эти classes не наследуются от общего base:
центральному workflow нужны только названные operations. Оба reporters должны
обработать 15 blocks и согласиться, что после этой работы complete первые 9.

\begin{shellcode}
cd examples/09_templates/generic_course_audit
cmake -S . -B build -DCMAKE_BUILD_TYPE=Debug
cmake --build build
./build/generic_course_audit --root ../../..
ctest --test-dir build --output-on-failure
\end{shellcode}

\begin{outputcode}
.........FFFFFF compact checked=15, passed=9
...
detailed checked=15, passed=9
reporters agree=true
GENERIC_COURSE_AUDIT=PASS
\end{outputcode}

Точки означают successful blocks, \texttt{F} --- будущие incomplete 10--15,
а не сбой программы. Дополнительный negative test передаёт
\texttt{BadReporter} без \texttt{on\_result}/\texttt{on\_finish}; compiler
отвергает соответствующую instantiation. В отличие от runtime Strategy блока~8,
reporter нельзя заменить неизвестным subtype во время run. Здесь это уместно:
набор output types известен при сборке и virtual hierarchy ничего не добавила
бы к предметному смыслу.

PASS блока требует: обе CMake-конфигурации C++17, warning-clean builds, 5
CTest cases, оба positive markers, две сохранённые negative diagnostics и
общую сборку PDF. Sanitizers и benchmark --- N/A: практика не содержит
memory-safety hypothesis или performance claim.

\subsection{STL начинается с этой идеи}

\texttt{std::vector<T>} --- class template, поэтому element type ---
частью concrete vector type. \texttt{std::array<T, N>} дополнительно получает
compile-time size. Generic algorithms описываются через требования к данным и
операциям, а не копируются для каждого element type.

Это только мост. Мы ещё не изучаем container taxonomy, iterator categories,
complexity contracts и standard algorithms: они составят блок~14. Теперь их
архитектура не будет выглядеть магией --- templates дают языковой фундамент
для независимости container/algorithm от одного concrete element type.

\subsection{Проверка понимания}

Нужно уметь объяснить:
1) какое дублирование устраняет function template; 2) различия template, specialization и instantiation; 3) что выводит deduction и почему смешанные arguments могут конфликтовать; 4) почему generic не означает «для любого типа»; 5) requirements \texttt{maximum}, \texttt{count\_matching} и \texttt{run\_audit}; 6) почему \texttt{Box<int>} и \texttt{Box<std::string>} --- разные types; 7) почему \texttt{N} в \texttt{FixedBuffer<T, N>} известен compile time; 8) зачем template definition обычно видима в месте instantiation; 9) когда выбрать template policy, а когда runtime Strategy; 10) compile-time, diagnostic и code-size costs; 11) как templates подводят к STL, не заменяя изучение её contracts.

\subsection{Источники для сверки}

Термины template parameters, deduction, implicit/explicit instantiation и
explicit specialization сверены с C++17 working draft N4659, clauses
[temp.param], [temp.deduct], [temp.inst] и [temp.expl.spec]. Базовая
педагогическая последовательность сопоставлена с MIT 6.096 Advanced I;
устаревший стиль источника заменён на C++17 и \texttt{std::}. Связь с
\texttt{vector}/\texttt{array} проверена по предыдущим блокам и
\textit{Modern C++ Tutorial}; performance claims ограничены моделью ISO C++
Performance TR и не выданы за гарантии generated code.
